\documentclass[12pt]{article}
\usepackage[utf8]{inputenc}
\usepackage[T1]{fontenc}
\usepackage{hyperref}
\title{Punan Batu Language Resources for K-12 Curriculum Development}
\author{Compiled by Claude Sonnet 5.5 (claude-sonnet-5-5) for Andr\'as Kornai}
\date{}
\begin{document}
\maketitle

\section{Introduction}
Punan Batu (ISO 639-3 pnm) is an Austronesian (North Borneo, Outer Central Sarawak) language of Sarawak, Malaysia. Sources describe it as severely endangered; the commonly cited figure is about 30 speakers (Wurm 2000), which is old. Speakers are reported near Long Geng, southeast of Belaga. Some know a sung ritual language, Latala or Menirak. Search found no educational materials.

\section{Primary Educational Resources}
None found in search results.

\section{Linguistic and Reference Materials}
\begin{itemize}
\item Punan Batu language, Wikipedia: \url{https://en.wikipedia.org/wiki/Punan_Batu_language}
\item Endangered Languages Project entry: \url{https://endangeredlanguages.com/lang/2527?hl=en}
\end{itemize}

\section{Religious and Community Resources}
\begin{itemize}
\item Joshua Project, Batu Punan in Malaysia (lists Malay as primary language, so of limited linguistic value): \url{https://joshuaproject.net/people_groups/10730/MY}
\item Languagehat post on the Punan Batu of Borneo: \url{https://languagehat.com/the-punan-batu-of-borneo/}
\end{itemize}

\section{Digital and Audio Resources}
No audio or software resources were found.

\section{Curriculum Development Context}
With so few speakers and no found teaching materials, any programme would be a documentation and revitalization effort led by the community, probably alongside Malay-medium schooling.
\end{document}